% GENERATED FILE. Edit canonical lessons, then run npm run generate:books.
% canonical-source-hash: fnv1a64:45e192b35f47bd6e
% canonical-lessons: UR-C183-mulaqat-karna, UR-C183-shadi-karna, UR-C183-jina, UR-C183-marna, UR-C183-barhana

\chapter{To pay a visit, To marry, To be alive, To die, To increase}
\label{ch:ur-to-pay-a-visit-to-marry-to-be-alive-to-die-to-increase}
\hlchaptermodality{183}

\begin{chapteropening}
\textbf{By the end of this chapter:} \emph{I can say to pay a visit, to marry, to be alive, to die and to increase.}

Complete the last lesson of chapter 183: I can say to pay a visit, to marry, to be alive, to die and to increase.
\end{chapteropening}

\section[\texorpdfstring{\ur{ملاقات} \ur{کرنا} (mulāqāt karnā)}{ملاقات کرنا (mulāqāt karnā)}]{\ur{ملاقات} \ur{کرنا} (mulāqāt karnā) — to pay a visit, to meet (formally)}

\label{lesson:UR-C183-mulaqat-karna}



\begin{quote}
Before the new one: say the Urdu for to stand up, then the Urdu for to wait.
\end{quote}

\subsection*{You'll want to know: \ur{ملاقات} \ur{کرنا}}
\textbf{\ur{ملاقات} \ur{کرنا}} — \emph{mulāqāt karnā} — \textquotedblleft{}to pay a visit, to meet (formally)\textquotedblright{}.

\begin{cousinweb}
\textbf{\ur{ملاقات} \ur{کرنا}} (\emph{mulāqāt karnā}) joins \emph{mulāqāt}, \textquotedblleft{}a meeting\textquotedblright{}, and \emph{karnā}, \textquotedblleft{}to do\textquotedblright{}. \emph{mulāqāt} is Arabic, \textquotedblleft{}a meeting face to face\textquotedblright{}, and reached Urdu through Persian (Lane, root \emph{l-q-y}); \emph{karnā} comes from Sanskrit \emph{karoti}, \textquotedblleft{}does\textquotedblright{} (CDIAL 2814).
\end{cousinweb}

\begin{grammarlens}[title={What you've built}]
One more word for this chapter.
\end{grammarlens}

\subsection*{Your turn}
\begin{itemize}
\raggedright
  \item \emph{Say it:} \emph{mulāqāt karnā}
  \item \emph{Say it:} \emph{mulāqāt karnā}, once more
  \item \emph{Say it:} say \emph{mulāqāt karnā}
  \item \emph{Recall:} say the Urdu for to stand up, then the Urdu for to wait, then say \emph{mulāqāt karnā} again
\end{itemize}

\begin{tcolorbox}[breakable,colback=teal!4,colframe=teal!35!black,title={Before you move on}]
What does \ur{ملاقات} \ur{کرنا} mean? (\textquotedblleft{}To pay a visit, to meet (formally)\textquotedblright{}.) Say it once more.
\end{tcolorbox}
\section[\texorpdfstring{\ur{شادی} \ur{کرنا} (shādī karnā)}{شادی کرنا (shādī karnā)}]{\ur{شادی} \ur{کرنا} (shādī karnā) — to marry}

\label{lesson:UR-C183-shadi-karna}



\begin{quote}
Before the new one: say the Urdu for to wait, then the Urdu for to pay a visit.
\end{quote}

\subsection*{You'll want to know: \ur{شادی} \ur{کرنا}}
\textbf{\ur{شادی} \ur{کرنا}} — \emph{shādī karnā} — \textquotedblleft{}to marry\textquotedblright{}.

\begin{cousinweb}
\textbf{\ur{شادی} \ur{کرنا}} (\emph{shādī karnā}) joins \emph{shādī}, \textquotedblleft{}a wedding\textquotedblright{}, and \emph{karnā}, \textquotedblleft{}to do\textquotedblright{}. \emph{shādī} is usually traced to Persian \emph{shādī}, \textquotedblleft{}joy\textquotedblright{}; \emph{karnā} comes from Sanskrit \emph{karoti}, \textquotedblleft{}does\textquotedblright{} (CDIAL 2814).
\end{cousinweb}

\begin{grammarlens}[title={What you've built}]
One more word for this chapter.
\end{grammarlens}

\subsection*{Your turn}
\begin{itemize}
\raggedright
  \item \emph{Say it:} \emph{shādī karnā}
  \item \emph{Say it:} \emph{shādī karnā}, once more
  \item \emph{Say it:} say \emph{shādī karnā}
  \item \emph{Recall:} say the Urdu for to wait, then the Urdu for to pay a visit, then say \emph{shādī karnā} again
\end{itemize}

\begin{tcolorbox}[breakable,colback=teal!4,colframe=teal!35!black,title={Before you move on}]
What does \ur{شادی} \ur{کرنا} mean? (\textquotedblleft{}To marry\textquotedblright{}.) Say it once more.
\end{tcolorbox}
\section[\texorpdfstring{\ur{جینا} (jīnā)}{جینا (jīnā)}]{\ur{جینا} (jīnā) — to be alive, to live (a life)}

\label{lesson:UR-C183-jina}



\begin{quote}
Before the new one: say the Urdu for to pay a visit, then the Urdu for to marry.
\end{quote}

\subsection*{You'll want to know: \ur{جینا}}
\textbf{\ur{جینا}} — \emph{jīnā} — \textquotedblleft{}to be alive, to live (a life)\textquotedblright{}.

\begin{cousinweb}
\textbf{\ur{جینا}} (\emph{jīnā}) is a tadbhava, a word inherited from Sanskrit and reshaped by sound change. It comes from \emph{jīvati}, \textquotedblleft{}is alive\textquotedblright{} (CDIAL 5241).
\end{cousinweb}

\begin{grammarlens}[title={What you've built}]
One more word for this chapter.
\end{grammarlens}

\subsection*{Your turn}
\begin{itemize}
\raggedright
  \item \emph{Say it:} \emph{jīnā}
  \item \emph{Say it:} \emph{jīnā}, once more
  \item \emph{Say it:} say \emph{jīnā}
  \item \emph{Recall:} say the Urdu for to pay a visit, then the Urdu for to marry, then say \emph{jīnā} again
\end{itemize}

\begin{tcolorbox}[breakable,colback=teal!4,colframe=teal!35!black,title={Before you move on}]
What does \ur{جینا} mean? (\textquotedblleft{}To be alive, to live (a life)\textquotedblright{}.) Say it once more.
\end{tcolorbox}
\section[\texorpdfstring{\ur{مرنا} (marnā)}{مرنا (marnā)}]{\ur{مرنا} (marnā) — to die}

\label{lesson:UR-C183-marna}



\begin{quote}
Before the new one: say the Urdu for to marry, then the Urdu for to be alive.
\end{quote}

\subsection*{You'll want to know: \ur{مرنا}}
\textbf{\ur{مرنا}} — \emph{marnā} — \textquotedblleft{}to die\textquotedblright{}.

\begin{cousinweb}
\textbf{\ur{مرنا}} (\emph{marnā}) is a tadbhava, a word inherited from Sanskrit and reshaped by sound change. It comes from \emph{marate}, \textquotedblleft{}will die\textquotedblright{} (CDIAL 9871). CDIAL calls direct descent from this old form possible, even probable.
\end{cousinweb}

\begin{grammarlens}[title={What you've built}]
One more word for this chapter.
\end{grammarlens}

\subsection*{Your turn}
\begin{itemize}
\raggedright
  \item \emph{Say it:} \emph{marnā}
  \item \emph{Say it:} \emph{marnā}, once more
  \item \emph{Say it:} say \emph{marnā}
  \item \emph{Recall:} say the Urdu for to marry, then the Urdu for to be alive, then say \emph{marnā} again
\end{itemize}

\begin{tcolorbox}[breakable,colback=teal!4,colframe=teal!35!black,title={Before you move on}]
What does \ur{مرنا} mean? (\textquotedblleft{}To die\textquotedblright{}.) Say it once more.
\end{tcolorbox}
\section[\texorpdfstring{\ur{بڑھانا} (baṛhānā)}{بڑھانا (baṛhānā)}]{\ur{بڑھانا} (baṛhānā) — to increase, to raise}

\label{lesson:UR-C183-barhana}



\begin{quote}
Before the new one: say the Urdu for to be alive, then the Urdu for to die.
\end{quote}

\subsection*{You'll want to know: \ur{بڑھانا}}
\textbf{\ur{بڑھانا}} — \emph{baṛhānā} — \textquotedblleft{}to increase, to raise\textquotedblright{}.

\begin{cousinweb}
\textbf{\ur{بڑھانا}} (\emph{baṛhānā}) is the causing form of \emph{baṛhnā}, \textquotedblleft{}to increase, to grow\textquotedblright{}, which CDIAL lists under Sanskrit \emph{vardhate}, \textquotedblleft{}grows, increases\textquotedblright{} (CDIAL 11376).
\end{cousinweb}

\begin{grammarlens}[title={What you've built}]
One more word for this chapter.
\end{grammarlens}

\subsection*{Your turn}
\begin{itemize}
\raggedright
  \item \emph{Say it:} \emph{baṛhānā}
  \item \emph{Say it:} \emph{baṛhānā}, once more
  \item \emph{Say it:} say \emph{baṛhānā}
  \item \emph{Recall:} say the Urdu for to be alive, then the Urdu for to die, then say \emph{baṛhānā} again
\end{itemize}

\begin{tcolorbox}[breakable,colback=teal!4,colframe=teal!35!black,title={Before you move on}]
What does \ur{بڑھانا} mean? (\textquotedblleft{}To increase, to raise\textquotedblright{}.) Say it once more.
\end{tcolorbox}
